\documentclass[a4paper,10pt]{article}
\usepackage[a4paper, top=2cm, bottom=1.8cm, left=1.5cm, right=1.5cm]{geometry}
\usepackage{amsmath, amssymb, amsfonts}
\usepackage{multicol}
\usepackage{enumitem}
\usepackage{xcolor}
\usepackage{fancyhdr}
\usepackage{mdframed}
\usepackage{array}
\usepackage{booktabs}
\usepackage{parskip}

% ─── Colors ──────────────────────────────────────────────────────────────────
\definecolor{darkblue}{RGB}{0,51,102}
\definecolor{lightgray}{RGB}{240,240,240}
\definecolor{accentred}{RGB}{180,0,0}

% ─── Page style ──────────────────────────────────────────────────────────────
\pagestyle{fancy}
\fancyhf{}
\setlength{\headheight}{14pt}
\renewcommand{\headrulewidth}{1.2pt}
\fancyhead[L]{\small\textcolor{darkblue}{\textbf{Department of Computer Engineering}}}
\fancyhead[R]{\small\textcolor{darkblue}{Page \thepage}}
\fancyfoot[C]{\small\textcolor{gray}{\textit{--- End of Page \thepage\ ---}}}

% ─── Part Section banner ─────────────────────────────────────────────────────
\newcommand{\partsection}[1]{%
  \vspace{5pt}%
  \noindent\colorbox{darkblue}{%
    \makebox[\dimexpr\columnwidth-2\fboxsep][c]{%
      \bfseries\color{white}\small #1%
    }%
  }%
  \vspace{4pt}\par%
}


% ─── Instruction box ─────────────────────────────────────────────────────────
\newmdenv[
  linecolor=darkblue,
  linewidth=1pt,
  backgroundcolor=lightgray,
  innertopmargin=5pt,
  innerbottommargin=5pt,
  innerleftmargin=7pt,
  innerrightmargin=7pt,
  skipabove=3pt,
  skipbelow=3pt
]{instructbox}

\setlength{\columnsep}{1cm}
\setlength{\parindent}{0pt}
\setlength{\parskip}{2pt}
\setlist[enumerate]{topsep=2pt, itemsep=1pt, parsep=0pt, leftmargin=1.5em}

% =============================================================================
\begin{document}
% =============================================================================

% ── HEADER BLOCK ─────────────────────────────────────────────────────────────
\begin{center}
    {\Large\bfseries\color{darkblue} MIDTERM EXAMINATION IN ADVANCED ENGINEERING MATHEMATICS}\\[4pt]
    {\small\color{black} Bachelor of Science in Computer Engineering \quad|\quad Academic Year 2025--2026}\\[6pt]
    \begin{tabular}{@{}p{0.46\textwidth} p{0.46\textwidth}@{}}
        \textbf{Name:}\enspace\underline{\hfill\hspace{6cm}}    &
        \textbf{Date:}\enspace\underline{\hfill\hspace{6cm}}      \\[6pt]
        \textbf{Section:}\enspace\underline{\hfill\hspace{6cm}} &
        \textbf{Score:}\enspace\underline{\hfill\hspace{6cm}}     \\
    \end{tabular}
\end{center}

\vspace{5pt}

\textbf{GENERAL INSTRUCTIONS:} This examination contains \textbf{20 items}. Show \textbf{all work} clearly and in a logical, step-by-step manner. Unsupported final answers will not receive full credit. Put your final answers in a \textbf{box}.

\vspace{4pt}
\begin{multicols}{2}

    %% ═══════════════════════════════════════════════════════════
    %%  PART I – SERIES
    %% ═══════════════════════════════════════════════════════════
    \noindent Items 1--2: Determine whether the given series converges or diverges. If it converges, find the \textbf{Radius of Convergence (RoC)} and \textbf{Interval of Convergence (IoC)}, testing endpoints explicitly.

    \medskip
    \noindent\textbf{1.}
    \[
        \sum_{n=1}^{\infty} \frac{(-1)^{n}(x-3)^{n}}{n\cdot 4^{n}}
    \]

    \medskip
    \noindent\textbf{2.}
    \[
        \sum_{n=0}^{\infty} \frac{n!\,(2x-1)^{n}}{(3n)!}
    \]

    \medskip

    \textbf{3.} Find the \textbf{Taylor series} of $f(x)=\sin x$ centered at $a=\dfrac{\pi}{4}$. Write the \textbf{first four nonzero terms} and the general term. Simplify all coefficients.

    \medskip

    \textbf{4.} Find the \textbf{Maclaurin series} of $f(x)=x^{2}e^{-3x}$. Write the \textbf{first four nonzero terms} and the general term of the series.

    \vspace{3pt}
    %% ═══════════════════════════════════════════════════════════
    %%  PART II – SERIES SOLUTIONS
    %% ═══════════════════════════════════════════════════════════
    \noindent Items 5--9 (4 pts each): Classify $x=0$ as an \textbf{Ordinary Point (OP)}, \textbf{Regular Singular Point (RSP)}, or \textbf{Irregular Singular Point (ISP)}. Justify rigorously using the analyticity of $xp(x)$ and $x^{2}q(x)$, or the analyticity of $P(x)$ and $Q(x)$ in the standard form $y''+P(x)y'+Q(x)y=0$.

    \medskip
    \noindent\textbf{5.}
    \[
        (x^{2}+4)\,y'' + 3x\,y' - y = 0
    \]

    \medskip
    \noindent\textbf{6.}
    \[
        x^{2}y'' + (x+x^{3})\,y' + (x-1)\,y = 0
    \]

    \medskip
    \noindent\textbf{7.}
    \[
        x^{3}y'' + \cos(x)\,y' + x\,y = 0
    \]

    \medskip
    \noindent\textbf{8.}
    \[
        (x-3)^{2}y'' + (x-3)\,y' + (x^{2}-9)\,y = 0
    \]

    \medskip
    \noindent\textbf{9.}
    \[
        2x\,y'' + y' + x\,y = 0
    \]

    \medskip

    \textbf{10.} Using the power series method $y=\displaystyle\sum_{n=0}^{\infty}a_{n}x^{n}$ about $x=0$:
    \[
        (1-x^{2})\,y'' - x\,y' + 4\,y = 0
    \]
    \begin{enumerate}[label=(\alph*)]
        \item Substitute into the ODE and derive the \textbf{recurrence relation}.
        \item Write the \textbf{first four nonzero terms} of each linearly independent solution $y_{1}$ and $y_{2}$ (with $y_1$ corresponding to $a_0\neq 0,\,a_1=0$ and $y_2$ to $a_0=0,\,a_1\neq 0$).
    \end{enumerate}


    \medskip
    \textbf{11.}
    \[
        4x\,y'' + 2\,y' + y = 0
    \]
    \begin{enumerate}[label=(\alph*)]
        \item Confirm $x=0$ is a \textbf{regular singular point}.
        \item Derive the \textbf{indicial equation} and solve for the indicial roots $r_1, r_2$.
        \item Using Frobenius series $y=\displaystyle\sum_{n=0}^{\infty}a_n x^{n+r}$, find the \textbf{first three nonzero terms} for each independent solution.
    \end{enumerate}

    \medskip
    \noindent\textbf{12.}
    \[
        x^{2}y'' - x\,y' + y = 0
    \]
    \begin{enumerate}[label=(\alph*)]
        \item Apply the Frobenius method to obtain the \textbf{indicial equation}. Identify the \textbf{repeated root}.
        \item Write the \textbf{first solution} $y_1$ from the Frobenius series.
        \item State the \textbf{second linearly independent solution} $y_2$ (use the reduction of order / logarithmic form).
        \item Verify both solutions using the substitution $y=x^m$ for the standard Cauchy--Euler method.
    \end{enumerate}


    \vspace{3pt}
    %% ═══════════════════════════════════════════════════════════
    %%  PART III – VECTORS
    %% ═══════════════════════════════════════════════════════════

    \textbf{13.} A particle moves from point $A(1,-2,3)$ to point $B(4,2,-1)$.
    \begin{enumerate}[label=(\alph*)]
        \item Find the \textbf{displacement vector} $\overrightarrow{AB}$.
        \item Compute the \textbf{magnitude} $\|\overrightarrow{AB}\|$.
        \item Find the \textbf{direction angle} $\alpha$ that $\overrightarrow{AB}$ makes with the positive $x$-axis (to the nearest tenth of a degree).
    \end{enumerate}

    \medskip

    \textbf{14.} Given $\mathbf{u}=\langle 3,-1,2\rangle$ and $\mathbf{v}=\langle -1,4,2\rangle$:
    \begin{enumerate}[label=(\alph*)]
        \item Compute $\mathbf{u}\cdot\mathbf{v}$.
        \item Find the \textbf{angle} $\theta$ between $\mathbf{u}$ and $\mathbf{v}$ (to the nearest tenth of a degree).
        \item Classify the relationship: \textit{orthogonal}, \textit{parallel}, or \textit{neither}. Justify.
    \end{enumerate}

    \medskip

    \textbf{15.} Given $\mathbf{u}=\langle 1,2,-1\rangle$ and $\mathbf{v}=\langle 3,-1,2\rangle$:
    \begin{enumerate}[label=(\alph*)]
        \item Compute $\mathbf{u}\times\mathbf{v}$ via the determinant expansion.
        \item Verify orthogonality: show $(\mathbf{u}\times\mathbf{v})\cdot\mathbf{u}=0$ and $(\mathbf{u}\times\mathbf{v})\cdot\mathbf{v}=0$.
        \item Find the \textbf{unit normal vector} $\hat{n}=\dfrac{\mathbf{u}\times\mathbf{v}}{\|\mathbf{u}\times\mathbf{v}\|}$.
    \end{enumerate}

    \medskip

    \textbf{16.} Apply the \textbf{Gram--Schmidt process} to the linearly independent set
    \[
        \mathbf{v}_1=\langle 1,1,0\rangle,\quad
        \mathbf{v}_2=\langle 1,0,1\rangle,\quad
        \mathbf{v}_3=\langle 0,1,1\rangle
    \]
    to produce an \textbf{orthonormal basis} $\{\hat{\mathbf{e}}_1, \hat{\mathbf{e}}_2, \hat{\mathbf{e}}_3\}$ for $\mathbb{R}^{3}$. Show all intermediate orthogonal vectors.

    \vspace{3pt}
    %% ═══════════════════════════════════════════════════════════
    %%  PART IV – MATRICES
    %% ═══════════════════════════════════════════════════════════

    \textbf{17.} Solve the system using \textbf{Gaussian elimination with back-substitution}. Write the augmented matrix and clearly label each row operation performed.
    \[
        \begin{cases}
            \phantom{-}x + 2y - z = \phantom{-}6 \\
            2x - y + 3z = -3                     \\
            -x + 3y - 2z = 14
        \end{cases}
    \]

    \medskip

    \textbf{18.} For the matrix below, find all \textbf{eigenvalues} $\lambda$ and their corresponding \textbf{eigenvectors}. Write the characteristic equation, factor it, and solve for each eigenspace.
    \[
        A = \begin{bmatrix} 4 & 1 & 0 \\ 2 & 3 & 0 \\ 0 & 0 & 5 \end{bmatrix}
    \]

    \medskip

    \textbf{19.} The ciphertext \textbf{``UNGJMUKA''} was encrypted using the \textbf{Hill Cipher} with key matrix
    \[
        K = \begin{bmatrix} 3 & 2 \\ 5 & 7 \end{bmatrix}, \quad A=0,\;B=1,\;\ldots,\;Z=25.
    \]
    \begin{enumerate}[label=(\alph*)]
        \item Compute $K^{-1}\pmod{26}$ using the adjugate method.
        \item Decrypt the ciphertext letter-pair by letter-pair, showing the full modular matrix multiplication for each block.
    \end{enumerate}

    \medskip

    \textbf{20.}  A message was sent and was encoded using \textbf{Hamming (7,4)}. The received codeword is $\mathbf{r}=[\,0,\;1,\;1,\;0,\;1,\;1,\;1\,]$.
    \begin{enumerate}[label=(\alph*)]
        \item Calculate each \textbf{syndrome check bit} $c_1$, $c_2$, $c_3$ from $\mathbf{r}$.
        \item Interpret the syndrome $[c_3\,c_2\,c_1]$ in binary to \textbf{locate the erroneous bit}.
        \item Write the \textbf{corrected codeword}.
    \end{enumerate}

\end{multicols}

\vspace{5pt}
\noindent\textcolor{darkblue}{\rule{\linewidth}{1.2pt}}
\begin{center}
    \small\textit{``The measure of intelligence is the ability to change.''} ---
    \textcolor{darkblue}{\textbf{--- END OF EXAMINATION ---}}
\end{center}


\vspace{20pt}
\noindent
\begin{tabular}{@{}>{\centering\arraybackslash}p{0.45\textwidth}
    @{\hspace{0.1\textwidth}}
    >{\centering\arraybackslash}p{0.45\textwidth}@{}}
    Prepared By:                             & Approved By:                                   \\[22pt]
    \underline{\textbf{KAIVENI TOM DAGCUTA}} & \underline{\textbf{ENGR. RODESITA S. ESTENZO}} \\
    \textit{Instructor}                      & \textit{Department Head}                       \\
\end{tabular}

\end{document}